% profiles/university-of-florida.tex — University of Florida
% ============================================================================
% source: https://grad.ufl.edu/current/academics/tdp/formatting/
%         (the Graduate School's current "Formatting Requirements & Templates"
%          page — the prose requirements)
% source: https://grad.ufl.edu/media/gradufledu/ms-word/TDP-Title-Page.docx
%         (the Graduate School's own title-page sample, linked from that page)
% source: https://it.ufl.edu/helpdesk/graduate-resources/ms-word--latex-templates/
%         (the Thesis & Dissertation Support Center's templates page, which the
%          Graduate School page mandates — see the delegation note below. The
%          accessibility requirements are stated here and nowhere else.)
% source: https://it.ufl.edu/helpdesk/media/helpdeskufledu/LaTeX_Example_Template.zip
%         (UF's OFFICIAL LaTeX template, `ufdissertation.cls` inside it. The
%          title page below is that class's \maketitle, and the accessibility
%          declarations below are its exampleMasterFile.tex \DocumentMetadata.)
% accessed: 2026-09-05
% checked:  texlib-thesis-profile-round
%
% VERIFIED and set here: geometry, spacing, the title page, the absence of an
% approval page, and three accessibility requirements. NOT set: chapter-opening
% margin (no such rule).
%
% ============================================================================
% READ THIS BEFORE FILING AT UF. THE GRADUATE SCHOOL MANDATES ITS OWN TEMPLATE.
%
%   "The UFIT Help Desk's Thesis & Dissertation Support Center maintains
%    templates for MS Word and LaTeX which you MUST use to format your thesis or
%    dissertation to make adherence to our formatting requirements easier."
%
% and, for LaTeX specifically:
%
%   "To ensure that your LaTeX document will be accessible you must use the
%    provided template and run your compiler with LuaLaTeX and TeXlive 2025."
%
% This profile makes texlib-thesis produce a document that MATCHES UF's rules.
% It does not make texlib-thesis into UF's template, and no profile can. UF is
% the first institution in this set to mandate a specific class file, and a
% filer who ignores that is taking a risk this file cannot remove. The profile
% is still worth having — for drafting, for checking a figure, and because the
% Editorial Office reviews the document rather than the source — but the
% mandate is real and is stated here rather than buried.
% ============================================================================
%
% ON THE DELEGATION. Under RESEARCHING.md only the graduate school's own
% requirements or its official template count as sources. The UFIT pages qualify
% on both counts and only because of that: the Graduate School's own formatting
% page names the UFIT Support Center's templates and says students must use
% them, so those templates ARE the graduate school's official template, and the
% accessibility rules attached to them are the rules a UF filing must meet.
% Nothing here is taken from UFIT as an independent authority.
%
% STALE LINK, so a later pass does not chase it: the UFIT page's "GRADUATE
% SCHOOL FORMATTING GUIDE" points at https://grad.ufl.edu/academics/editorial/etd-specs/,
% which is a 404. The Graduate School reorganised its site; the live path is
% /current/academics/tdp/formatting/. Several other plausible URLs
% (grad.ufl.edu/academics/tdp/, success.grad.ufl.edu/td/formatting) are also
% dead or redirect to the homepage.

\thesisinstitution{University of Florida}

% --- Filing figures ----------------------------------------------------------
% "Margins: One inch (1") all around (top, right, bottom, left) on all pages."
% Stated exactly, with no range and no exception — note "on all pages", which
% rules out a different margin on the title page or on chapter openings.
% Independently confirmed twice: the Word title-page sample's section properties
% are 1440 twips on all four sides, and UF's own LaTeX class declares
% \geometry{margin=1in,paperheight=11in,paperwidth=8.5in}.
\thesissetgeometry{left=1in, right=1in, top=1in, bottom=1in}

% "Spacing: Double-space paragraph text. Indent the first line of each
% paragraph. Single-space headings, table titles, and figure captions. The
% reference list is single-spaced with one blank line in between the
% references."
%
% Double is REQUIRED here for body text, not offered as one of several options —
% the first UF-style rule in this round that is a rule rather than a choice. The
% single-spaced elements are per-element and not expressible in a profile.
\thesissetspacing{double}

% No chapter-opening rule. "One inch all around on all pages" positively
% forbids a deeper one, so \thesissetchapteropening stays unset — here that is a
% verified absence, not an unknown.

% --- Accessibility -----------------------------------------------------------
% ALL THREE CHECK-FLAGS ARE RAISED, and the evidence is unusually direct.
%
% The prose mandate: "Keep in mind that your document must meet accessibility
% standards", and "Due to new accessibility requirements, you must use LuaLaTex
% and Texlive 2025 in your compiler." Modal "must", twice, on the page the
% Graduate School's own formatting page sends students to.
%
% The specific values come from UF's OFFICIAL LaTeX template rather than from
% prose: exampleMasterFile.tex opens with
%
%   \DocumentMetadata{lang=en-US, pdfstandard=ua-2, pdfversion=2.0,
%                     tagging=on, tagging-setup={math/setup=mathml-SE}}
%
% That is the school declaring, in the file it requires students to use, that a
% filed UF dissertation is a tagged PDF/UA-2 document with its language set. A
% reviewer who thinks a template's own metadata is weaker evidence than prose
% should say so — this is the one place in this profile where a requirement is
% read off a template rather than out of a sentence, and it is flagged for that
% reason.
\thesisrequiretagging
\thesisrequiredocumentlanguage
\thesisrequirepdfstandard{PDF/UA-2}
\thesisaccessibilityrequirement{Accessible document}
	{"Keep in mind that your document must meet accessibility standards."
	 UF names no external standard such as WCAG in the filing requirements; what
	 it names is the workflow below and the tagged output it produces.}
\thesisaccessibilityrequirement{LaTeX workflow}
	{"To ensure that your LaTeX document will be accessible you must use the
	 provided template and run your compiler with LuaLaTeX and TeXlive 2025."
	 This class satisfies the engine half of that and NOT the template half —
	 see the header of this file.}
\thesisaccessibilityrequirement{Alt text and tables}
	{UF's own testing notes that alt text on images and TikZ diagrams, MathML
	 mathematics and tables with header rows all work in this workflow, and
	 that "Table headers have to be declared prior to and external to a table,
	 as LaTeX has no native table header element" — i.e. set table/header-rows
	 in \string\DocumentMetadata rather than expecting the table to say so.}
\thesisaccessibilityrequirement{Known gaps, from UF}
	{UF publishes what does NOT work in this workflow: "Figure captions create
	 problematic tag structures", "Numbered equations also create problematic
	 tag structures", page headers and footers fall outside the reading order,
	 and "Document metadata, such as Title, does not populate properly into PDF
	 metadata". Recorded because a filer will otherwise think these are their
	 own errors.}

% --- Title page --------------------------------------------------------------
% This is UF's own \maketitle from ufdissertation.cls, reproduced. The Word
% title-page sample carries the same six blocks in the same order, so the two
% official sources agree.
%
% Everything is centred and uppercase except the article "A" and "By", which the
% class sets as written. Specifics carried over:
%
%   * The title is uppercased (\UF@TitleCase defaults to \MakeUppercase) and set
%     single-spaced, in a document that is otherwise double-spaced.
%   * The author is uppercased.
%   * The four-line statement is broken exactly as UF breaks it and must not be
%     reflowed.
%   * The date line is the degree YEAR ALONE — UF's class prints \degreeYear and
%     never \degreeMonth here. So a document filing at UF should set
%     \graddate{2027}, not \graddate{May, 2027}. This is the one place where the
%     class's single free-text date field does not match UF's two, and it is
%     resolved by instruction rather than by parsing the string.
%   * The page carries no number.
%
% UF's class opens with \vspace*{-0.4in} to lift the title above the 1in top
% margin. That is not reproduced: the requirement is "one inch all around on ALL
% pages", the negative space contradicts it, and a profile should follow the
% stated rule over a template's convenience. Noted so the difference is not
% mistaken for an oversight.
\renewcommand{\thesistitlepage}{%
	\loadgeometry{nopagenumber}%
	\begin{titlepage}%
		\thispagestyle{empty}\singlespacing\centering
		\thesis@tagtitle{Title}{\MakeUppercase{\@title}}\par
		\vfill
		By\\[\baselineskip]
		\MakeUppercase{\@author}\par
		\vfill
		A \MakeUppercase{\thesis@DocNoun} PRESENTED TO THE GRADUATE SCHOOL\\
		OF THE UNIVERSITY OF FLORIDA IN PARTIAL FULFILLMENT\\
		OF THE REQUIREMENTS FOR THE DEGREE OF\\
		\MakeUppercase{\thesis@degree}\\[\baselineskip]
		UNIVERSITY OF FLORIDA\\[\baselineskip]
		\thesis@date\par
	\end{titlepage}%
	\loadgeometry{normal}%
}

% --- Committee approval page --------------------------------------------------
% UF's filed document has none. The Graduate School's page order is exhaustive
% and reads: "Title Page, Copyright Page, Dedication (if any), Acknowledgements,
% Table of Contents, List of Tables (if any), List of Figures (if any), List of
% Objects (if any), List of Abbreviations (if any), Abstract, Chapters,
% Appendices (if any), List of References, Biographical Sketch." No approval,
% committee or signature page appears in it, and UF's own LaTeX class emits none.
%
% PHYSICAL SIGNATURES: none are collected in the document. UF's formatting
% requirements and template prescribe no signature line anywhere, and the
% supervisory chair is metadata in the template (\chair{...}) rather than a page.
% Approval is handled through the Graduate School's submission workflow, outside
% the manuscript.
\thesisnoapprovalpage
